%!TEX root = ../../main.tex
% ============================================================
% 几何作图 / 点子图与观察物体
% 本文件由 main.tex 通过 \input 引入，请勿单独编译
% 重构日期: 2026-04-11
% ============================================================

\subsection{解答：在点子图上画面积相等的图形}

\vspace{0.4cm}

\noindent\textbf{1. 在点子图上画和长方形面积相等的平行四边形、三角形和梯形各一个。}

\vspace{1.2cm}

\begin{center}
% =====================================================================
% 【整体绘制思路：等面积图形的构造】
% 1. 点子图背景：
%    使用 \foreach 循环绘制 21x7 的点阵阵列，模拟题目要求的“点子图”环境。
%    点阵原点为 (0,0)，最大至 (20,6)。
% 2. 图形面积逻辑运算 (目标面积为 8)：
%    - 给定长方形：宽 2，高 4。计算面积 = 2 * 4 = 8。
% 3. 构造等积作图图形 (所有高统一取最直观的 4 单位)：
%    - 平行四边形：面积 = 底 * 高。需要 底 * 4 = 8，得出 底 = 2。
%    - 三角形：面积 = 底 * 高 / 2。需要 底 * 4 / 2 = 8，得出 底 = 4。
%    - 梯形：面积 = (上底 + 下底) * 高 / 2。需要 (上底 + 下底) * 4 / 2 = 8，
%            得出 (上底 + 下底) = 4。可以取 上底=1，下底=3。
% 4. 坐标排布规则：
%    - 每个图形之间保持 1~2 个单位的间隔，均匀分布在点阵图上。
%    - 为了清晰，已知图形用黑色表示，要求画出的图形全部用红色表现。
% =====================================================================
\begin{tikzpicture}[>=Stealth, line width=1.0pt, scale=0.8]

  % =========== 背景点子图 ===========
  % \foreach \x ... \foreach \y ... 批量画出所有的点
  \foreach \x in {0,1,2,3,4,5,6,7,8,9,10,11,12,13,14,15,16,17,18,19,20} {
    \foreach \y in {0,1,2,3,4,5,6} {
      \fill[gray!80] (\x,\y) circle (1.8pt);
    }
  }

  % =========== 原图形：长方形 (黑色) ===========
  % 宽为2，长(高)为4，面积 = 2 * 4 = 8
  % 左下角定于 (1, 1)
  \draw[black, thick] (1,1) rectangle (3,5);

  % =========== 解答图形1：平行四边形 (红色) ===========
  % 底为2，高为4，面积 = 2 * 4 = 8
  % 从 x=4 开始画底，底从 (4,1) 到 (6,1)
  % 上底整体向右平移 1 个单位，从 (5,5) 到 (7,5)
  \draw[red!70!black, line width=1.3pt] (4,1) -- (6,1) -- (7,5) -- (5,5) -- cycle;

  % =========== 解答图形2：三角形 (红色) ===========
  % 底为4，高为4，面积 = 4 * 4 / 2 = 8
  % 从 x=8 开始画底，底从 (8,1) 到 (12,1)
  % 顶点选取中间位置确保视觉对称，即 (10,5)
  \draw[red!70!black, line width=1.3pt] (8,1) -- (12,1) -- (10,5) -- cycle;

  % =========== 解答图形3：梯形 (红色) ===========
  % 上底1，下底3，高为4，面积 = (1 + 3) * 4 / 2 = 8
  % 从 x=14 开始画下底，下底从 (14,1) 到 (17,1)
  % 参考原答案图，左侧边垂直（直角梯形），所以上底从 (14,5) 开始
  % 上底长1，到 (15,5)
  \draw[red!70!black, line width=1.3pt] (14,1) -- (17,1) -- (15,5) -- (14,5) -- cycle;

\end{tikzpicture}
\end{center}

\vspace{0.8cm}

{\small\color{gray}
\textbf{解析：}\\
长方形的面积 = 长 $\times$ 宽 = $2 \times 4 = 8$。\\
题目要求画出的图形面积都需要为 $8$。为方便在点阵图作图，我们统一规定这三个图形的高都取 $4$（占用4个竖直单元格）。\\
1. \textbf{平行四边形的面积} = 底 $\times$ 高。已知面积=$8$，高=$4$，所以**底必须为 $2$**。\\
2. \textbf{三角形的面积} = 底 $\times$ 高 $\div\ 2$。已知面积=$8$，高=$4$，可列式：底 $\times 4 \div 2 = 8$，解得**底必须为 $4$**。\\
3. \textbf{梯形的面积} = (上底 + 下底) $\times$ 高 $\div\ 2$。已知面积=$8$，高=$4$，可列式：(上底 + 下底) $\times 4 \div 2 = 8$，解得 (上底 + 下底)= $4$。结合点阵特点，可以取**上底为 $1$，下底为 $3$**（直角梯形画法最清晰）。
}

\vspace{0.6cm}

{\small\color{blue!70!black}
\textbf{绘图基本原理与步骤：}\\
\textbf{1. 点子图阵列构建}：利用 TikZ 强大的 \texttt{\char`\\foreach} 嵌套循环功能，快速生成一整片均匀排列的圆点（\texttt{circle (1.8pt)}），作为辅助坐标系。\\
\textbf{2. 整数格点捕捉}：所有的图形顶点（如 \texttt{(4,1)}, \texttt{(7,5)}）都被精确设定在整数坐标上，确保线段端点能够地落在中国点子阵列图的几何视觉点上，严格遵循了“在点子图上作图”的要求。\\
\textbf{3. 色彩区分层级}：原图矩形选用黑色刻画题目条件，要求绘制的平行四边形、三角形、梯形统一选用视觉鲜明的深红色（\texttt{red!70!black}，线宽 \texttt{1.3pt}），突出作图结果。
}

\newpage

\newpage

\subsection{解答：从上面看各是什么形状}

\vspace{0.4cm}

\noindent\textbf{1. 从上面看圆柱、正方体、圆锥和长方体，看到的各是什么形状？连一连。}

\vspace{1.2cm}

\begin{center}
% =====================================================================
% 【整体绘制思路：立体图形俯视图连线】
% 1. 坐标系与布局：
%    - 上排 (立体图形)：y=5 水平线。从左到右依次为 圆柱(x=2)、正方体(x=6)、圆锥(x=10)、长方体(x=15)。
%    - 下排 (平面图形)：y=1 水平线。从左到右依次为 长方形(x=2)、带圆心的圆(x=6)、圆(x=10)、正方形(x=15)。
% 2. 三维模拟绘画法：
%    - 圆柱：上下两个椭圆，两侧垂直连线。
%    - 正方体/长方体：绘制正面矩形，再通过固定偏移向量绘制顶面和侧面，形成等大透视感。
%    - 圆锥：底面画完整椭圆和虚线后半段，连接顶点。
% 3. 连线锚点与线型：
%    - 为每个图形的正下方(T)和正上方(B)预埋 \coordinate 锚点。
%    - 依据正确答案：圆柱连圆、正方体连正方形、圆锥连带点圆、长方体连长方形，使用红色粗线绘制。
% =====================================================================
\begin{tikzpicture}[>=Stealth, line width=1.0pt, scale=0.85]

  % =========== 上排：立体图形 ===========
  
  % 1. 圆柱 (Cylinder)
  \begin{scope}[shift={(2, 5)}]
    \draw (0, 0.8) ellipse (0.9 and 0.25);
    \draw (-0.9, 0.8) -- (-0.9, -0.6);
    \draw (0.9, 0.8) -- (0.9, -0.6);
    \draw (-0.9, -0.6) arc (180:360:0.9 and 0.25);
    % 预留连接锚点
    \coordinate (T1) at (0, -1.0);
  \end{scope}

  % 2. 正方体 (Cube)
  \begin{scope}[shift={(6, 5)}]
    % 正面 (尺寸 1.4 x 1.4)
    \draw (-0.7, -0.7) rectangle (0.7, 0.7);
    % 利用偏移量 (0.5, 0.5) 画顶面和右侧面
    \coordinate (FrontTL) at (-0.7, 0.7);
    \coordinate (FrontTR) at (0.7, 0.7);
    \coordinate (FrontBR) at (0.7, -0.7);
    \coordinate (BackTL)  at (-0.2, 1.2);
    \coordinate (BackTR)  at (1.2, 1.2);
    \coordinate (BackBR)  at (1.2, -0.2);
    
    \draw (FrontTL) -- (BackTL) -- (BackTR) -- (FrontTR);
    \draw (FrontTR) -- (BackTR) -- (BackBR) -- (FrontBR);
    \coordinate (T2) at (0, -1.0);
  \end{scope}

  % 3. 圆锥 (Cone)
  \begin{scope}[shift={(10, 5)}]
    % 底部椭圆弧
    \draw (-1.0, -0.6) arc (180:360:1.0 and 0.25);
    \draw[dashed] (1.0, -0.6) arc (0:180:1.0 and 0.25);
    % 连线到顶点
    \draw (-1.0, -0.6) -- (0, 1.3) -- (1.0, -0.6);
    \coordinate (T3) at (0, -1.0);
  \end{scope}

  % 4. 长方体 (Rectangular Cuboid)
  \begin{scope}[shift={(15, 5)}]
    % 正面 (尺寸 2.4 x 1.0，宽而扁)
    \draw (-1.2, -0.5) rectangle (1.2, 0.5);
    % 利用偏移量 (0.6, 0.5) 画顶面和右侧面
    \coordinate (FrontTL4) at (-1.2, 0.5);
    \coordinate (FrontTR4) at (1.2, 0.5);
    \coordinate (FrontBR4) at (1.2, -0.5);
    \coordinate (BackTL4)  at (-0.6, 1.0);
    \coordinate (BackTR4)  at (1.8, 1.0);
    \coordinate (BackBR4)  at (1.8, 0);
    
    \draw (FrontTL4) -- (BackTL4) -- (BackTR4) -- (FrontTR4);
    \draw (FrontTR4) -- (BackTR4) -- (BackBR4) -- (FrontBR4);
    \coordinate (T4) at (0.3, -1.0); % 因为向右延展，连线锚点稍微右移找中心
  \end{scope}


  % =========== 下排：平面视图 ===========

  % A. 长方形
  \begin{scope}[shift={(2, 1)}]
    \draw (-1.1, -0.6) rectangle (1.1, 0.6);
    \coordinate (B1) at (0, 1.2);
  \end{scope}

  % B. 带圆心的圆
  \begin{scope}[shift={(6, 1)}]
    \draw (0, 0) circle (0.9);
    \fill (0, 0) circle (1.5pt);
    \coordinate (B2) at (0, 1.2);
  \end{scope}

  % C. 普通圆
  \begin{scope}[shift={(10, 1)}]
    \draw (0, 0) circle (0.9);
    \coordinate (B3) at (0, 1.2);
  \end{scope}

  % D. 正方形
  \begin{scope}[shift={(15, 1)}]
    \draw (-0.7, -0.7) rectangle (0.7, 0.7);
    \coordinate (B4) at (0, 1.2);
  \end{scope}

  % =========== 绘制答案连线 (红色) ===========
  % \draw[color, line width] (起点) -- (终点);
  
  % 1. 圆柱从上面看是：普通圆 (C)
  \draw[red!80!black, line width=1.5pt] (T1) -- (B3);
  
  % 2. 正方体从上面看是：正方形 (D)
  \draw[red!80!black, line width=1.5pt] (T2) -- (B4);
  
  % 3. 圆锥从上面看是：带圆心的圆 (B)
  \draw[red!80!black, line width=1.5pt] (T3) -- (B2);
  
  % 4. 长方体从上面看是：长方形 (A)
  \draw[red!80!black, line width=1.5pt] (T4) -- (B1);

\end{tikzpicture}
\end{center}

\vspace{0.8cm}

{\small\color{gray}
\textbf{解析：}\\
观察立体图形的三视图（尤其是俯视图），想象从物体正上方往下看所能看到的轮廓平面图。\\
1. \textbf{圆柱}：上下底面都是圆，侧面垂直，所以从正上方往下看，只能看到一个与底面相同的**圆**。\\
2. \textbf{正方体}：所有的六个面都是完全相等的正方形，从正上方往下看，看到的就是其顶面，是一个**正方形**。\\
3. \textbf{圆锥}：底面是一个圆，最上面是一个尖顶。从正上方俯视时，能看到底面的圆形轮廓，并且能看到正中央那个尖顶（在平面图上体现为一个圆心点）。因此俯视图是**一个带有一点（圆心）的圆**。\\
4. \textbf{长方体}：相对的两个面是全等的长方形，从正上方往下看，看到的是它原本作为顶面的那个**长方形**。
}

\vspace{0.6cm}

{\small\color{blue!70!black}
\textbf{绘图基本原理与步骤：}\\
\textbf{1. 伪 3D 效果手工绘制}：纯利采用二维基础命令线构建极具透视感的各种立体图。比如正方形与长方体利用平行四边形偏移法画出顶面与侧边；圆柱和圆锥利用 \texttt{ellipse} (椭圆) 来呈现圆面在三维视角下的透视形变，圆锥隐蔽的背面底弧则贴心地使用了虚线 \texttt{dashed}。\\
\textbf{2. 预设解耦描点连线结构}：图形绘制采用相互独立、定距排列的 \texttt{scope} 分组模块化编写。并在每个模块的图形下方/上方埋下了不可见的透明锚点（如 \texttt{T1,T2} 以及 \texttt{B1,B2}）。当绘制连线时，仅需将对应的锚点对接即可，使得相交乱序的红线在代码层面上依然非常清爽清晰，且极易修正错题。\\
\textbf{3. 色彩重点剥离}：黑色的图作为客观存在的认知题干选项；红色加粗连线（\texttt{red!80!black}）强引力充当核心作答答案，对比鲜明符合辅导教材的美学要求。
}


\newpage
